\documentclass[10pt,a4paper]{amsart}
\usepackage[margin=19mm]{geometry}
\usepackage{amsmath,amssymb,mathtools}
\usepackage[scr=rsfs,cal=euler]{mathalfa}
\usepackage{xfrac}
\usepackage[unicode,hidelinks,bookmarks=false]{hyperref}
\newcommand{\ie}{i.e.\ }
\newcommand{\cf}{cf.\ }
\newcommand{\resp}{respectively\ }
\newcommand{\Subsection}{\subsection}
\providecommand{\nobreakdash}{\nobreak\hbox{-}}
\setlength{\emergencystretch}{2em}
\multlinegap0pt
\usepackage{tikz}
\usepackage{mathtools}
\usepackage{amssymb}
\newtheorem{prop}{Proposition}
\newtheorem{theorem}{Theorem}
\usepackage{cleveref}
\crefname{prop}{Proposition}{Propositions}
\crefname{theorem}{Theorem}{Theorems}
\newcommand{\R}{\mathbb{R}}
\title{An explicit construction of Kaleidocycles by elliptic theta functions}
\author{Shizuo Kaji, Kenji Kajiwara, Shota Shigetomi}
\date{Source: arXiv:2308.04977v4 (CC BY 4.0)}
\begin{document}
\maketitle

\noindent\emph{Source and licence.} An explicit construction of Kaleidocycles by elliptic theta functions. Version arXiv:2308.04977v4 (CC BY 4.0), distributed by arXiv under the Creative Commons Attribution 4.0 International licence, http://creativecommons.org/licenses/by/4.0/, as recorded in the arXiv licence metadata for this exact version and shown on the abstract page https://arxiv.org/abs/2308.04977v4. Full licence notice: \texttt{licenses/arxiv-2308.04977-CC-BY-4.0.txt}.\par\smallskip

\noindent\textit{Editorial scope.} Counted unit: the article's theorem thm:PhDKal2 (source0.tex lines 647--683). The article's complete original proof of it is delivered with it (source0.tex lines 684--834, one \texttt{proof} environment of 151 lines). No mathematics is written, repaired or re-derived: the statement and the proof are the article's own lines, in the article's own notation, and no retained line is substituted or re-flowed.  Retained uncounted as the source-local context the proof needs: lines 389--396; lines 400--438; lines 576--616; lines 593--595; lines 2305--2324; lines 2331--2338.  Nothing was omitted from any retained span, so the package carries no omission record.  No retained line contains a citation, so the wrapper carries no bibliography and no bibliography entry.  No retained line contains a figure environment, an image-inclusion command or drawing code, so no image is delivered and the package declares no asset dependency.  Editorial formatting only, in addition: an AMS \texttt{amsart} preamble that restates the article's own declarations and macros used by the retained lines, namely the environments \texttt{prop}, \texttt{theorem} on separate counters, together with the standard packages those lines need.  The retained environments print in this document as Proposition 1, Theorem 1 in order of appearance, whereas the article prints the same items with the numbering of its own full document.  The complete native source of the article as distributed in the arXiv e-print for this exact version is bundled unchanged as \texttt{sources/145418/source0.tex}.
\begin{align}
&f_{n}f_{n}^*+g_{n}g_{n}^*=F_{n+1}F_{n} \label{i1},\\
&\beta_{1} F_{n}f_{n}+G_{n}g_{n}^*=F_{n+1}f_{n-1} \label{i2},\\
&\beta_{1} F_{n}g_{n}-G_{n}f_{n}^*=F_{n+1}g_{n-1} \label{i3},\\
&H_{n+1}F_{n}-H_{n}F_{n+1}=f_{n}^*g_{n} \label{i4},\\
&D_zF_{n+1}\cdot F_{n}:=\frac{dF_{n+1}}{dz}F_{n}-F_{n+1}\frac{dF_{n}}{dz}=g_{n}g_{n}^* \label{i5},\\
&\left|G_{n}\right|^{2}=\widetilde{G}_{n}^{2},\label{i6}
\end{align}

\begin{align}\label{T}
  T_{n} &=\left. \frac{1}{f_{n}f_{n}^{*}+g_{n}g_{n}^{*}} \left(
    \begin{array}{c}
   \vspace{2mm}
      f_{n}^{*}g_{n}+f_{n}g_{n}^{*} \\
      \vspace{2mm}
      \frac{1}{i}\left(f_{n}^{*}g_{n}-f_{n}g_{n}^{*}\right) \\
      f_{n}f_{n}^{*}-g_{n}g_{n}^{*}
    \end{array}
  \right)\right|_{z=0}, \\
  \widetilde{N}_{n}&=\left.\frac{1}{f_{n}f_{n}^{*}+g_{n}g_{n}^{*}} \left(
    \begin{array}{c}
    \vspace{2mm}
      \left\{(f_{n}^{*})^{2}-(g_{n}^{*})^{2}\right\}\varrho_{n}+ \left\{(f_{n})^{2}-(g_{n})^{2}\right\}\varrho_{n}^{*} \\
      \vspace{2mm}
      \frac{1}{i}\left[\left\{(f_{n}^{*})^{2}+(g_{n}^{*})^{2}\right\}\varrho_{n}-\left\{(f_{n})^{2}+(g_{n})^{2}\right\}\varrho_{n}^{*}\right] \\
      -2\left(f_{n}^{*}g_{n}^{*}\varrho_{n}+f_{n}g_{n}\varrho_{n}^{*}\right)
    \end{array}
  \right)\right|_{z=0}, \nonumber\\
  \widetilde{B}_{n}&=\left. \frac{1}{f_{n}f_{n}^{*}+g_{n}g_{n}^{*}} \left(
    \begin{array}{c}
    \vspace{2mm}
       \frac{1}{i}\left[-\left\{(f_{n}^{*})^{2}-(g_{n}^{*})^{2}\right\}\varrho_{n}+ \left\{(f_{n})^{2}-(g_{n})^{2}\right\}\varrho_{n}^{*}\right] \\
       \vspace{2mm}
      \left\{(f_{n}^{*})^{2}+(g_{n}^{*})^{2}\right\}\varrho_{n}+ \left\{(f_{n})^{2}+(g_{n})^{2}\right\}\varrho_{n}^{*} \\
      \frac{2}{i}\left(f_{n}^{*}g_{n}^{*}\varrho_{n}-f_{n}g_{n}\varrho_{n}^{*}\right)
    \end{array}
  \right)\right|_{z=0}, \nonumber \\
  %
 \gamma_{n}&= \left.\left(
        \begin{array}{c}
        \vspace{2mm}
    \dfrac{H_{n}+H_{n}^*}{F_{n}}\\
    \vspace{2mm}
    \dfrac{1}{i}\dfrac{H_{n}-H_{n}^*}{F_{n}}\\
    n-2\dfrac{\partial\log F_{n}}{\partial z}
        \end{array}
      \right) \right|_{z=0}. \nonumber  
\end{align}

\begin{prop}\label{thm:PhDKal1}
Put $\mu_{n}=ivn+\frac{z}{\Delta_{3}-\Delta_{1}}$
and define
\begin{align}\label{t1}
&F_{n}(z)=\alpha_1\exp\left(\frac{n\Delta_{3}}{\Delta_{3}-\Delta_{1}}z\right)\vartheta_2\left(\mu_{n}-\frac{1}{2}iv\right),\\
&f_{n}(z)=\alpha_2 R_{3}^{n}\exp\left(\frac{\left(n+\frac{1}{2}\right)\Delta_{3}}{\Delta_{3}-\Delta_{1}}z\right)\vartheta_2\left(\mu_{n}+r\right), \nonumber\\
&g_{n}(z)=\alpha_3 R_{1}^{-n}\exp\left(\frac{\left(n+\frac{1}{2}\right)\Delta_{3}}{\Delta_{3}-\Delta_{1}}z\right)\vartheta_4\left(\mu_{n}-r\right), \nonumber\\
&G_{n}(z)=\alpha_4 R_{1}^{-n}R_{3}^{n}\exp\left(\frac{n\Delta_{3}}{\Delta_{3}-\Delta_{1}}z\right)\vartheta_4\left(\mu_{n}-\frac{1}{2}iv\right), \nonumber\\
&H_{n}(z)=\alpha_5 R_{1}^{-n}R_{3}^{-n}\exp\left(\frac{n\Delta_{3}}{\Delta_{3}-\Delta_{1}}z\right)\vartheta_4\left(\mu_{n}-\frac{1}{2}iv-2r\right), \nonumber \\
&\widetilde{G}_{n}(z)=-i\alpha_4\exp\left(\frac{n\Delta_{3}}{\Delta_{3}-\Delta_{1}}z\right)\vartheta_4\left(\mu_{n}-\frac{1}{2}iv\right), \nonumber
\end{align}
where
\begin{align}
&\alpha_1=\sqrt{\vartheta_3(2r)\vartheta_3(0)},\quad \alpha_2=\vartheta_3\left(-\frac{1}{2}iv+r\right),\quad \alpha_3=\vartheta_1\left(\frac{1}{2}iv+r\right),\nonumber\\
&\alpha_4=\frac{\vartheta_1(iv)\sqrt{\vartheta_3(2r)\vartheta_3(0)}}{\vartheta_3(0)},\quad \alpha_5=\frac{u_{1}}{\vartheta_3(2r)\vartheta_1(iv)\sqrt{\vartheta_3(2r)\vartheta_3(0)}}.\nonumber
\end{align}
Then, they satisfy \eqref{i1}-\eqref{i6} with
\begin{equation}\label{beta_1}
\beta_{1}=\frac{\vartheta_3(iv)}{\vartheta_3(0)},
\end{equation}
and form $\tau$ functions.
The corresponding discrete curve $\gamma$ defined by \eqref{T} has a unit segment length $\left|\gamma_{n+1}-\gamma_{n}\right|=1$.
The torsion angle is given by
\begin{equation}\label{lam1}
\begin{split}
\lambda_n&=\arccos\left(\frac{G_{n}^{*}(0)G_{n-1}(0)+G_{n}(0)G_{n-1}(0)^{*}}{2\widetilde{G}_{n}(0)\widetilde{G}_{n-1}(0)}\right)\\
&=\arccos\left(\frac{R_{1}R_{3}^{-1}+R_{1}^{-1}R_{3}}{2}\right).
\end{split}
\end{equation}
In particular, this does not depend on $n$,
and $\gamma$ has a constant torsion angle
$\lambda_n=\lambda_0$.
Moreover,
the signed curvature angle is given by
\begin{equation}\label{kappa1}
\begin{split}
\kappa_{n}&=2\arctan\left(\frac{\widetilde{G}_{n}(0)}{\beta_{1} F_{n}(0)}\right)\\
&=2\arctan\left(-i\frac{\vartheta_1(iv)\vartheta_4\left(ivn-\frac{1}{2}iv\right)}{\vartheta_3(iv)\vartheta_2\left(ivn-\frac{1}{2}iv\right)}\right).
\end{split}
\end{equation}
\end{prop}

\begin{equation}
\beta_{1}=\frac{\vartheta_3(iv)}{\vartheta_3(0)},
\end{equation}

\begin{align}
&\vartheta_2(X_1)\vartheta_2(Y_1)\vartheta_3(U_1)\vartheta_3(V_1)-\vartheta_4(X_1)\vartheta_4(Y_1)\vartheta_1(U_1)\vartheta_1(V_1)\label{pr4-1}\\
&=\vartheta_2(X)\vartheta_2(Y)\vartheta_3(U)\vartheta_3(V)-\vartheta_4(X)\vartheta_4(Y)\vartheta_1(U)\vartheta_1(V),\nonumber\\
&\vartheta_4(X_1)\vartheta_3(Y_1)\vartheta_1(U_1)\vartheta_2(V_1)-\vartheta_2(X_1)\vartheta_1(Y_1)\vartheta_3(U_1)\vartheta_4(V_1)\label{pr4-2}\\
&=\vartheta_4(X)\vartheta_3(Y)\vartheta_1(U)\vartheta_2(V)-\vartheta_2(X)\vartheta_1(Y)\vartheta_3(U)\vartheta_4(V),\nonumber\\
&\vartheta_2(X_1)\vartheta_2(Y_1)\vartheta_4(U_1)\vartheta_4(V_1)+\vartheta_1(X_1)\vartheta_1(Y_1)\vartheta_3(U_1)\vartheta_3(V_1)\label{pr4-3}\\
&=\vartheta_4(X)\vartheta_4(Y)\vartheta_2(U)\vartheta_2(V)-\vartheta_3(X)\vartheta_3(Y)\vartheta_1(U)\vartheta_1(V),\nonumber\\
&\vartheta_3(X_1)\vartheta_3(Y_1)\vartheta_2(U_1)\vartheta_2(V_1)+\vartheta_1(X_1)\vartheta_1(Y_1)\vartheta_4(U_1)\vartheta_4(V_1)\label{pr4-3-1}\\
&=\vartheta_3(X)\vartheta_3(Y)\vartheta_2(U)\vartheta_2(V)+\vartheta_1(X)\vartheta_1(Y)\vartheta_4(U)\vartheta_4(V),\nonumber\\
&\vartheta_4(X_1)\vartheta_4(Y_1)\vartheta_4(U_1)\vartheta_4(V_1)-\vartheta_1(X_1)\vartheta_1(Y_1)\vartheta_1(U_1)\vartheta_1(V_1)\label{pr4-3-2}\\
&=\vartheta_4(X)\vartheta_4(Y)\vartheta_4(U)\vartheta_4(V)-\vartheta_1(X)\vartheta_1(Y)\vartheta_1(U)\vartheta_1(V),\nonumber\\
&\vartheta_4(X_1)\vartheta_4(Y_1)\vartheta_1(U_1)\vartheta_1(V_1)-\vartheta_1(X_1)\vartheta_1(Y_1)\vartheta_4(U_1)\vartheta_4(V_1)\label{pr4-3-3}\\
&=\vartheta_4(X)\vartheta_4(Y)\vartheta_1(U)\vartheta_1(V)-\vartheta_1(X)\vartheta_1(Y)\vartheta_4(U)\vartheta_4(V),\nonumber\\
&\vartheta_4(X_1)\vartheta_4(Y_1)\vartheta_1(U_1)\vartheta_1(V_1)+\vartheta_2(X_1)\vartheta_2(Y_1)\vartheta_3(U_1)\vartheta_3(V_1)\label{pr4-3-4}\\
&=\vartheta_3(X)\vartheta_3(Y)\vartheta_2(U)\vartheta_2(V)-\vartheta_1(X)\vartheta_1(Y)\vartheta_4(U)\vartheta_4(V),\nonumber\\
&\vartheta_3(X_1)\vartheta_3(Y_1)\vartheta_1(U_1)\vartheta_1(V_1)-\vartheta_1(X_1)\vartheta_1(Y_1)\vartheta_3(U_1)\vartheta_3(V_1)\label{pr4-3-6}\\
&=\vartheta_3(X)\vartheta_3(Y)\vartheta_1(U)\vartheta_1(V)-\vartheta_1(X)\vartheta_1(Y)\vartheta_3(U)\vartheta_3(V),\nonumber\\
&\vartheta_3(X_1)\vartheta_3(Y_1)\vartheta_4(U_1)\vartheta_4(V_1)-\vartheta_1(X_1)\vartheta_1(Y_1)\vartheta_2(U_1)\vartheta_2(V_1)\label{pr4-3-7}\\
&=\vartheta_3(X)\vartheta_3(Y)\vartheta_4(U)\vartheta_4(V)-\vartheta_1(X)\vartheta_1(Y)\vartheta_2(U)\vartheta_2(V).\nonumber
\end{align}

\begin{equation}\label{cl}
\begin{split}
&\vartheta_1\left(X+niy; iy\right)=(-1)^{n}\exp\{-\pi i\left(2nX+n^{2}iy\right)\}\vartheta_1\left(X;iy\right),\\
&\vartheta_2\left(X+niy;iy\right)=\exp\{-\pi i\left(2nX+n^{2}iy\right)\}\vartheta_2\left(X;iy\right),\\
&\vartheta_3\left(X+niy;iy\right)=\exp\{-\pi i\left(2nX+n^{2}iy\right)\}\vartheta_3\left(X;iy\right),\\
&\vartheta_4\left(X+niy;iy\right)=(-1)^{n}\exp\{-\pi i\left(2nX+n^{2}iy\right)\}\vartheta_4\left(X;iy\right).
\end{split}
\end{equation}

\begin{theorem}\label{thm:PhDKal2}
Denote $\mu_{n}=ivn+\frac{z}{\Delta_{3}-\Delta_{1}}$. Consider the following functions parameterised by $t\in \R$:
\begin{align}\label{t}
&F_{n}(t, z)=\alpha_1\exp\left(\frac{n\Delta_{3}}{\Delta_{3}-\Delta_{1}}z+\frac{C}{2}tz\right)\vartheta_2\left(\mu_{n}-\frac{1}{2}iv+it\right),\\
&f_{n}(t, z)=\alpha_2 R_{3}^{n}\exp\left(\frac{\left(n+\frac{1}{2}\right)\Delta_{3}}{\Delta_{3}-\Delta_{1}}z+\frac{C}{2}tz-\frac{\Gamma}{2} it\right)\vartheta_2\left(\mu_{n}+r+it\right), \nonumber\\
&g_{n}(t, z)=\alpha_3 R_{1}^{-n}\exp\left(\frac{\left(n+\frac{1}{2}\right)\Delta_{3}}{\Delta_{3}-\Delta_{1}}z+\frac{C}{2}tz+\frac{\Gamma}{2} it\right)\vartheta_4\left(\mu_{n}-r+it\right), \nonumber\\
&G_{n}(t, z)=\alpha_4 R_{1}^{-n}R_{3}^{n}\exp\left(\frac{n\Delta_{3}}{\Delta_{3}-\Delta_{1}}z+\frac{C}{2}tz\right)\vartheta_4\left(\mu_{n}-\frac{1}{2}iv+it\right), \nonumber\\
&H_{n}(t, z)=\alpha_5 R_{1}^{-n}R_{3}^{-n}\exp\left(\frac{n\Delta_{3}}{\Delta_{3}-\Delta_{1}}z+\frac{C}{2}tz+\Gamma it\right)\vartheta_4\left(\mu_{n}-\frac{1}{2}iv-2r+it\right), \nonumber\\
&\widetilde{G}_{n}(t, z)=-i\alpha_4 \exp\left(\frac{n\Delta_{3}}{\Delta_{3}-\Delta_{1}}z+\frac{C}{2}tz\right)\vartheta_4\left(\mu_{n}-\frac{1}{2}iv+it\right),\nonumber
\end{align}
where $C, \Gamma\in\mathbb{R}$ are constants.
Then, they satisfy \eqref{i1}-\eqref{i6} with $\beta_{1}$ given by \eqref{beta_1} and form $\tau$ functions.
For each $t$,
the corresponding discrete curve $\gamma$ \eqref{T} has a unit segment length $\left|\gamma_{n+1}-\gamma_{n}\right|=1$.
The torsion angle is given by
\begin{equation}\label{lam2}
\begin{split}
\lambda_n&=\arccos\left(\frac{G_{n}^{*}(t, 0)G_{n-1}(t, 0)+G_{n}(t, 0)G_{n-1}^{*}(t, 0)}{2\widetilde{G}_{n}(t, 0)\widetilde{G}_{n-1}(t, 0)}\right)\\
&=\arccos\left(\frac{R_{1}R_{3}^{-1}+R_{1}^{-1}R_{3}}{2}\right).
\end{split}
\end{equation}
In particular, this depends neither on $n$ nor on $t$,
and $\gamma$ has a constant torsion angle
$\lambda_n(t)=\lambda_0(0)$.
Moreover,
the signed curvature angle is given by
\begin{equation}\label{kappa2}
\begin{split}
\kappa_{n}\left(t\right)&=\kappa_{n}\left(t;v,y\right)=2\arctan\left(\frac{\widetilde{G}_{n}(t, 0)}{\beta_{1} F_{n}(t, 0)}\right)\\
&=2\arctan\left(-i\frac{\vartheta_1(iv)\vartheta_4\left(ivn-\frac{1}{2}iv+it\right)}{\vartheta_3(iv)\vartheta_2\left(ivn-\frac{1}{2}iv+it\right)}\right),
\end{split}
\end{equation}
and has a period in $t$:
\begin{equation}\label{kappa_periodic}
    \kappa_{n}\left(t+y\right)=-\kappa_{n}\left(t\right).
\end{equation}
\end{theorem}

\begin{proof}
Putting
\begin{eqnarray}\label{pr6-1}
\left\{
\begin{array}{l}
\vspace{2mm}
X=ivn+r+it+\frac{z}{\Delta_{3}-\Delta_{1}},\\
\vspace{2mm}
Y=ivn-r+it+\frac{z}{\Delta_{3}-\Delta_{1}},\\
\vspace{2mm}
U=\frac{1}{2}iv+r,\\
V=\frac{1}{2}iv-r,\\
\end{array}
\right.
\end{eqnarray}
in \eqref{pr4-1} yields \eqref{i1}.
Putting
\begin{eqnarray}\label{pr6-2}
\left\{
\begin{array}{l}
\vspace{2mm}
X=iv\left(n-\frac{1}{2}\right)+it+\frac{z}{\Delta_{3}-\Delta_{1}},\\
\vspace{2mm}
Y=ivn+r+it+\frac{z}{\Delta_{3}-\Delta_{1}},\\
\vspace{2mm}
U=iv,\\
V=\frac{1}{2}iv-r,\\
\end{array}
\right.
\end{eqnarray}
in \eqref{pr4-1} yields \eqref{i2}.
Putting
\begin{eqnarray}\label{pr6-3}
\left\{
\begin{array}{l}
\vspace{2mm}
X=iv\left(n-\frac{1}{2}\right)+it+\frac{z}{\Delta_{3}-\Delta_{1}},\\
\vspace{2mm}
Y=\frac{1}{2}iv+r,\\
\vspace{2mm}
U=iv,\\
V=ivn-r+it+\frac{z}{\Delta_{3}-\Delta_{1}},\\
\end{array}
\right.
\end{eqnarray}
in \eqref{pr4-2} yields \eqref{i3}.
Putting
\begin{eqnarray}\label{pr6-4}
\left\{
\begin{array}{l}
\vspace{2mm}
X=ivn-r+it+\frac{z}{\Delta_{3}-\Delta_{1}},\\
\vspace{2mm}
Y=-\frac{1}{2}iv+r,\\
\vspace{2mm}
U=ivn-r+it+\frac{z}{\Delta_{3}-\Delta_{1}},\\
V=-\frac{1}{2}iv+r,\\
\end{array}
\right.
\end{eqnarray}
in \eqref{pr4-3} yields
\begin{equation}\label{pr6-5}
\begin{split}
&\frac{\vartheta_2\left(0\right)\vartheta_4\left(0\right)}{\vartheta_1\left(-\frac{1}{2}iv+r\right)\vartheta_3\left(-\frac{1}{2}iv+r\right)}\vartheta_2\left(\mu_{n}-\frac{1}{2}iv+it\right)\vartheta_4\left(\mu_{n+1}-\frac{1}{2}iv-2r+it\right)\\
&=\frac{\vartheta_2\left(-\frac{1}{2}iv+r\right)\vartheta_4\left(-\frac{1}{2}iv+r\right)}{\vartheta_1\left(-\frac{1}{2}iv+r\right)\vartheta_3\left(-\frac{1}{2}iv+r\right)}\vartheta_2\left(\mu_{n}-r+it\right)\vartheta_4\left(\mu_{n}-r+it\right)\\
&-\vartheta_3\left(\mu_{n}-r+it\right)\vartheta_1\left(\mu_{n}-r+it\right).
\end{split}
\end{equation}
Also, putting
\begin{eqnarray}\label{pr6-6}
\left\{
\begin{array}{l}
\vspace{2mm}
X=ivn-r+it+\frac{z}{\Delta_{3}-\Delta_{1}},\\
\vspace{2mm}
Y=\frac{1}{2}iv+r,\\
\vspace{2mm}
U=ivn-r+it+\frac{z}{\Delta_{3}-\Delta_{1}},\\
V=\frac{1}{2}iv+r,\\
\end{array}
\right.
\end{eqnarray}
in \eqref{pr4-3} yields
\begin{equation}\label{pr6-7}
\begin{split}
&\frac{\vartheta_2\left(0\right)\vartheta_4\left(0\right)}{\vartheta_1\left(\frac{1}{2}iv+r\right)\vartheta_3\left(\frac{1}{2}iv+r\right)}\vartheta_2\left(\mu_{n+1}-\frac{1}{2}iv+it\right)\vartheta_4\left(\mu_{n}-\frac{1}{2}iv-2r+it\right)\\
&=\frac{\vartheta_2\left(\frac{1}{2}iv+r\right)\vartheta_4\left(\frac{1}{2}iv+r\right)}{\vartheta_1\left(\frac{1}{2}iv+r\right)\vartheta_3\left(\frac{1}{2}iv+r\right)}\vartheta_2\left(\mu_{n}-r+it\right)\vartheta_4\left(\mu_{n}-r+it\right)\\
&-\vartheta_3\left(\mu_{n}-r+it\right)\vartheta_1\left(\mu_{n}-r+it\right).
\end{split}
\end{equation}
Subtracting \eqref{pr6-7} from \eqref{pr6-5} and substituting 
\begin{equation}\label{pr6-8}
\begin{split}
&\vartheta_3\left(2r\right)\vartheta_4\left(0\right)\vartheta_2\left(0\right)\vartheta_1\left(iv\right)\\
&=\vartheta_1\left(\frac{1}{2}iv+r\right)\vartheta_3\left(\frac{1}{2}iv+r\right)\vartheta_2\left(-\frac{1}{2}iv+r\right)\vartheta_4\left(-\frac{1}{2}iv+r\right)\\
&-\vartheta_1\left(-\frac{1}{2}iv+r\right)\vartheta_3\left(-\frac{1}{2}iv+r\right)\vartheta_2\left(\frac{1}{2}iv+r\right)\vartheta_4\left(\frac{1}{2}iv+r\right),
\end{split}
\end{equation}
 in the equation, we obtain
\begin{equation}\label{p4-4}
\begin{split}
&\frac{\vartheta_1\left(-\frac{1}{2}iv+r\right)\vartheta_3\left(-\frac{1}{2}iv+r\right)\vartheta_1\left(\frac{1}{2}iv+r\right)\vartheta_3\left(\frac{1}{2}iv+r\right)}{\vartheta_3(2r)\vartheta_1(iv)}\\
&\times\left\{\frac{\vartheta_1\left(\frac{1}{2}iv+r\right)\vartheta_3\left(\frac{1}{2}iv+r\right)}{\vartheta_1\left(-\frac{1}{2}iv+r\right)\vartheta_3\left(-\frac{1}{2}iv+r\right)}\vartheta_2\left(\mu_{n}-\frac{1}{2}iv+it\right)\vartheta_4\left(\mu_{n+1}-\frac{1}{2}iv-2r+it\right)\right.\\
&\left.-\vartheta_2\left(\mu_{n+1}-\frac{1}{2}iv+it\right)\vartheta_4\left(\mu_{n}-\frac{1}{2}iv-2r+it\right)\right\}\\
&=\vartheta_3\left(\frac{1}{2}iv+r\right)\vartheta_2\left(\mu_{n}-r+it\right)\vartheta_1\left(\frac{1}{2}iv+r\right)\vartheta_4\left(\mu_{n}-r+it\right),
\end{split}
\end{equation}
from which we see \eqref{i4} holds.
Putting
\begin{eqnarray}\label{pr6-9}
\left\{
\begin{array}{l}
\vspace{2mm}
X=ivn+r+it+\frac{z}{\Delta_{3}-\Delta_{1}},\\
\vspace{2mm}
Y=ivn-r+it+\frac{z}{\Delta_{3}-\Delta_{1}},\\
\vspace{2mm}
U=\frac{1}{2}iv+r+h,\\
V=\frac{1}{2}iv-r+h,\\
\end{array}
\right.
\end{eqnarray}
in \eqref{pr4-1} yields
\begin{equation}\label{pr6-10}
\begin{split}
&\vartheta_2\left(\mu_{n+1}-\frac{1}{2}iv+it+h\right)\vartheta_2\left(\mu_{n}-\frac{1}{2}iv+it-h\right)\vartheta_3(2r)\vartheta_3(0)\\
&=\vartheta_{2}\left(\mu_{n}+r+it\right)\vartheta_{2}\left(\mu_{n}-r+it\right)\vartheta_{3}\left(\frac{1}{2}iv+r+h\right)\vartheta_{3}\left(\frac{1}{2}iv-r+h\right)\\
&-\vartheta_{4}\left(\mu_{n}+r+it\right)\vartheta_{4}\left(\mu_{n}-r+it\right)\vartheta_{1}\left(\frac{1}{2}iv+r+h\right)\vartheta_{1}\left(\frac{1}{2}iv-r+h\right).
\end{split}
\end{equation}
Differentiating \eqref{pr6-10} with $h$ and then putting $h=0$, we obtain
\begin{equation}\label{pr6-11}
\begin{split}
&\vartheta_3(2r)\vartheta_3(0)\left(\Delta_{3}-\Delta_{1}\right)D_{z}\vartheta_2\left(\mu_{n+1}-\frac{1}{2}iv+it\right)\cdot\vartheta_2\left(\mu_{n}-\frac{1}{2}iv+it\right)\\
&=-\Delta_{3}\vartheta_{3}\left(-\frac{1}{2}iv+r\right)\vartheta_{3}\left(\frac{1}{2}iv+r\right)\vartheta_{2}\left(\mu_{n}+r+it\right)\vartheta_{2}\left(\mu_{n}-r+it\right)\\
&-\Delta_{1}\vartheta_{1}\left(-\frac{1}{2}iv+r\right)\vartheta_{1}\left(\frac{1}{2}iv+r\right)\vartheta_{4}\left(\mu_{n}+r+it\right)\vartheta_{4}\left(\mu_{n}-r+it\right).
\end{split}
\end{equation}
Using the equation \eqref{pr6-10} with $h=0$ and substituting \eqref{pr6-11}, we obtain
\begin{equation}\label{p4-5}
\begin{split}
&\vartheta_3(2r)\vartheta_3(0)\left\{\left(\Delta_{3}-\Delta_{1}\right)D_{z}\vartheta_2\left(\mu_{n+1}-\frac{1}{2}iv+it\right)\cdot\vartheta_2\left(\mu_{n}-\frac{1}{2}iv+it\right)\right.\\
&\left.+\Delta_{3}\vartheta_2\left(\mu_{n+1}-\frac{1}{2}iv+it\right)\vartheta_2\left(\mu_{n}-\frac{1}{2}iv+it\right)\right\}\\
&=\left(\Delta_{3}-\Delta_{1}\right)\vartheta_1\left(-\frac{1}{2}iv+r\right)\vartheta_4\left(\mu_{n}+r+it\right)\vartheta_1\left(\frac{1}{2}iv+r\right)\vartheta_4\left(\mu_{n}-r+it\right).
\end{split}
\end{equation}
From this we see \eqref{i5} holds.
By the definition of $G$ and $\widetilde{G}$, \eqref{i6} follows.
Periodicity of the signed curvature angle \eqref{kappa_periodic} follows from \eqref{cl}.
\Cref{thm:PhDKal1} follows from Theorem \ref{thm:PhDKal2} by setting $t=0$.
\end{proof}

\end{document}
